\chapter{Human Direction as a Program}
\label{ch:32}

The argument of this book is that the dominant opposition between salvation and extinction poses the wrong question, and that a better one is available. The better question asks what humanity should request of systems of greater capability if such systems exist, and it is answered by specifying the structure within which any request is made. The structure has been built in six parts. The first fixed the terms: the problems of the present are real whether or not any new capability arrives, they are coupled in ways that can make a collection of manageable crises unmanageable, the capacity to correct them can itself be eroded by them, and they present different kinds of obstacle, of which only the scientific and engineering kinds respond directly to capability. The second part examined the instruments, and found that the generation of proposals is cheap and its verification dear, that the acceleration of any field is limited by the slower of the two, that models fail in ways that optimization against them exploits, and that verification is worth only as much as the independence of the verifier from the generator and the margin by which a guarantee survives the gap between model and world. The third part considered institutions and held that what can be delegated in education, care, and public service is bounded below by a core of attention and judgment that cannot be, and that this core, because it does not shrink with automation, becomes the binding constraint and the proper object of public investment. The fourth and fifth parts applied a common standard of specification to the programs of very large scale, and in doing so separated the programs that survive their own arithmetic from those that do not. The sixth examined the prospect of extreme longevity as a problem of biology, of personal continuity, of reliability, and of institutions, and then fixed the constraint that prevents the distant future from being used to override the present.

Three findings of the middle chapters deserve to be recalled, because they show the method at work. The first is that energy is rarely the binding constraint on the megascale programs, and is more often a distraction. The energy to lift a cubic meter of water a kilometer, or to bring a kilogram to orbit, is a few kilowatt-hours, and the difficulties lie elsewhere: in the fifth power of the diameter of a pipe, in the dynamic pressure at the exit of an accelerator, in the strength of a fiber that does not yet exist. The second is that several of the proposals fail on analysis, and that the book has said so. A compression tower to the edge of space is excluded by a scaling law that no choice of joint can alter, a rigid shell around a star by its instability and mass, a rigid ring near the ground by a demand for compressive strength greater than that of any known fiber in tension. Stating these failures is not an attack on the ambitions that prompted the proposals. It is a condition of taking them seriously, for the failures point to what must be different, in a material or in a sequence, for the ambition to be realized. The third is that the surviving programs share a feature: each depends, for its legitimacy as for its physics, on something outside the engineering. The freshwater corridors depend on the consent of the states they cross, the forest supports on the consent of the communities who live where they would stand, the mass accelerators on a regime of verification that does not exist, and the swarm of collectors on an authorization that no state can give by itself.

That last feature is the book's conclusion in small. The question which institutions may authorize a given program has no technical answer, and the book has not pretended otherwise. It has stated what an adequate answer must contain: an authorization that is multilateral where the effects are, that provides for the participation of those who are not parties, that allows a party to decline participation without loss of its existing rights, that is revocable by a procedure agreed in advance, and that is accompanied by a plan for termination that does not itself produce a harm worse than the one it ends. It has insisted that the accountability for every consequence rest on a person or an institution that can be named, questioned, and sanctioned, and not on an artifact, and in this it follows the interview that prompted the work. The interview's insight, that the attribution of agency to a machine relocates responsibility from those who can answer for it to an entity that cannot, is the premise from which the architecture of authority proceeds, and the book's contribution is to carry that premise through the engineering: to show where, in a program of robotic embodiment, orbital construction, or planetary restoration, the decisions are made that cannot be delegated, and to design the program so that they remain with persons.

Limitations should be stated plainly. The calculations of the book are of the first test only: they show that a concept is or is not consistent with its own premises, and they do not show that it is wise, safe, or affordable. Several of the concepts treated, the volsorial and intervolsorial pediments, the geothermal mass accelerators, the nested expandable structures, are described in the author's earlier notes in terms that leave their specifications incomplete, and the definitions adopted here are readings of those descriptions, to be superseded as the specifications are fixed. The numerical examples use stipulated parameters where parameters are not established, and each such place has been marked. The treatment of longevity assumes simple models of mortality and of memory that capture orders of magnitude and not the detail of any real population. The ordering of the readiness classes is a judgment, and a reasonable reader may reclassify a program on grounds that the book has not considered. The book has not attempted to settle the empirical questions on which the public debate turns, whether capable systems will arrive, how soon, and with what dispositions, and it does not need to, for its argument is conditional. If they do arrive, the programs described here, with their gates, their verifiers, and their floors, are one account of what to ask of them and how to keep the asking in human hands. If they do not, the problems remain, the instruments of the second part are still those with which they must be addressed, and the institutions of the third part are still those that must be built.

\section*{Formal Summary}

Collecting the constructions above, a \emph{human-directed program} is a tuple\[(\Omega,A,R_c,\Sigma,\kappa,\mathcal{D},\{\varphi_j\},r,w_{\min})\]in which $\Omega$ is the space of trajectories, $A$ the admissible set fixed by legitimate procedure and contained in $A_{\mathrm{deleg}}$, $R_c$ the realizable set at capability $c$, $\Sigma$ a selection procedure, $\kappa$ the classification of tasks into local autonomy, supervised execution, and nondelegable judgment, $\mathcal{D}$ the set of locally autonomous behaviors certified as robustly invariant, $\{\varphi_j\}$ a family of independent verification predicates, $r$ a readiness function, and $w_{\min}$ the floor of present obligations. The program is \emph{admissible} when (i) the realized trajectory is $\omega^\ast=\Sigma(R_c\cap A)$, where $\Sigma$ exercises meaningful direction in the sense of Chapter 13 and every duty arising from its decisions is covered, that is, $\kappa_{\mathrm{cov}}(h;\beta_h)=1$ for every harm event $h$; (ii) every deployed proposal lies in $\Pi_c^\ast$ and has been checked by tests whose errors are not nested with those of the proposer; (iii) no judgment task is performed by a machine and every locally autonomous behavior lies in $\mathcal{D}$; (iv) the plan is readiness-respecting, which is a screening condition accompanied by a separate analysis of implementation risk; and (v) the selected policy belongs to $\mathcal{A}_{\mathrm{adm}}$. The evaluative thesis is that the merit of capability is assessed first by strict enlargement of $R_c\cap A$ in the sense of inclusion, and by a measure $\mu$ only where its choice has been justified and its influence reported. The share $a(\omega)$ of tasks transferred to machines is not a proxy for admissibility, since admissibility depends on which tasks are transferred and not on how many, and the two comparisons are therefore not treated as measuring the same quantity.
